\documentclass[10pt,a4paper]{article}
\usepackage[margin=18mm]{geometry}
\usepackage[T2A]{fontenc}
\usepackage[utf8]{inputenc}
\usepackage[english,russian]{babel}
\usepackage{amsmath,amssymb,booktabs,graphicx,array}
\usepackage[unicode,colorlinks=true,urlcolor=blue,linkcolor=blue]{hyperref}
\setlength{\parindent}{0pt}
\setlength{\parskip}{5pt}
\emergencystretch=3em
\newcommand{\R}{\operatorname{R}}
\title{t-Flow и PFGM++ на Exp / Spiral:\\полные координаты и полный U-Net}
\author{}
\date{11 сентября 2026}
\begin{document}
\maketitle
\vspace{-12mm}

\section*{Вопрос и устройство сравнения}
Проверяем, правильно ли добавлены два негауссовских метода и достаточно ли
хорошо они обучаются для оценки LID. Для каждой модели заново обучаются
Exp и Spiral в двух исходных представлениях: все 30 координат и все 784 пикселя
канонического изображения $28\times28$. Истинный LID равен 2 для Exp и 1 для
Spiral. Всего восемь обучений, один seed на модель/задачу/представление.

Ни в обучении, ни в оценке LID нет подгонки ковариационной оболочки, уменьшения
числа координат или готового ответа в нормальных направлениях. Коэффициенты
используют общий spectral residual core ширины 512 с четырьмя блоками,
3005214 параметров. Изображения используют общий U-Net ширины 32,
892353 параметра. В каждом представлении архитектура и число параметров
совпадают у обоих методов. Разные архитектуры для двух представлений разрешены;
название PCA у исходного набора не меняет маршрут полного изображения.

Использованы неизменённые канонические массивы, 99000 fit и 1000 source-train
holdout объектов, все 1000 test объектов. Только fit определяет вектор среднего
и один общий RMS-множитель. Одинаковы 128000 обновлений, batch 256, AdamW
с шагом $2\cdot10^{-4}$, weight decay $10^{-6}$, clipping 1, пары противоположных
шумов и последние 8000 шагов cosine/EMA. Веса выбирает минимальный нативный
holdout loss среди общих raw/EMA кандидатов. Обучение не останавливается раньше.
Бюджет обновлений и число сетевых примеров совпадают; равенство FLOPs не заявлено.

\section*{Что именно обучается и измеряется}
Обозначим $X$ чистую нормированную точку, $Y=X+\lambda Z$ её зашумлённое
наблюдение, а $b(Y,\lambda)$ --- предсказанный posterior mean чистой точки.
$\lambda$ означает RMS шума на одну ambient-координату. У t-Flow $Z$ имеет
радиальное Student-t ядро с $\nu=5$, у PFGM++ дополнительная размерность 128,
то есть Student-t $\nu=128$. Один chi-square знаменатель общий для всех
координат объекта: $Z=G\sqrt{(\nu-2)/S}$, где $G\sim\mathcal N(0,I_N)$,
$S\sim\chi^2_\nu$ независимы. Так каждая координата имеет дисперсию 1.
Хвосты фиксированы для всех задач; шум не обрезается.
Это два нативных метода одной семьи ядер, а не две несвязанные семьи.

t-Flow сохраняет равномерное нативное время и noise MSE, PFGM++ --- нативный
EDM weighted denoising loss и lognormal sigma. Их sampler и loss различаются
по определению метода. У Student нативный scale равен
$\sigma=\lambda\sqrt{(\nu-2)/\nu}$; у PFGM++ радиус
$r=\lambda\sqrt{\nu-2}$. Нативное время t-Flow равно
$t=k/(k+\lambda)$, где $k=\sqrt{\nu/(\nu-2)}$, а состояние потока равно $tY$.
Lognormal параметры PFGM++ переведены в общие
fit-RMS единицы: mean $-1.2+\log2$, std 1.2, sigma-data 1; поддержка условно
ограничена общей сеткой через inverse CDF. Это не kernel-only причинный эксперимент.

Оценка LID --- скаляр $\R(y,\lambda)=\operatorname{tr}D_y b(y,\lambda)$.
Для нативной скорости t-Flow это $Nt+t(1-t)\operatorname{div}v$, для PFGM++ ---
$N-r\operatorname{div}u$, где $N$ --- полная размерность входа.
Дивергенция скорости $v$ берётся по нативному состоянию $tY$, а дивергенция
радиального поля $u$ --- по $Y$ при фиксированном $r$. Gaussian
добавка квадрата displacement к этим негауссовским ядрам не применяется.
Векторный benchmark дифференцирует все 30 координат точно; image benchmark
использует 64 полных Rademacher-пробы Hutchinson. Отдельная проверка на 32
фиксированных holdout точках вычисляет точный trace по всем 784 координатам.

\section*{Почему выбранный MAE сам по себе недостаточен}
Один масштаб выбирается минимумом $\operatorname{mean}_i|\R_i-d_i|$ на 1000
source-train holdout объектах, по общим 29 кандидатам от $1/256$ до 64.
После фиксации весов и масштаба вычисляется test. Основной pointwise Kneedle
использует те же 29 кандидатов без LID-меток. Границы и fallback сохраняются.
Этот канонический test уже использовался в предыдущих экспериментах;
нет утверждения о новом слепом подтверждающем наборе. Рецепт текущих восьми
обучений зафиксирован до их запуска, без подбора параметров по их test-ответам.

Контроль без выученной нелинейной поправки даёт
$b_0(y,\lambda)=y/(1+\lambda^2)$ и $\R_0=N/(1+\lambda^2)$.
Известный постоянный $d$ позволяет подобрать к нему масштаб даже без изучения
многообразия. В таблице ниже отдельный $\lambda_0$ выбран тем же holdout-правилом;
MAE одинаков на всех 1000 test точках, поскольку ответ постоянен.
\input{results/zero_head_table.tex}
Все четыре контроля проходят ранее заданный практический порог MAE$<0.5$.
Поэтому ниже он остаётся только описательной проверкой выбранного score.
Качество posterior оценивается независимо по точному закону и свежему шуму.

\section*{Результаты обученных моделей}
В таблице MAE --- среднее $|\R-d|$ по всем 1000 test точкам; меньше лучше.
95\% интервалы получены query bootstrap при фиксированных весах, масштабе и
реализации probe-шума. Они не описывают неопределённость повторного обучения.
Kneedle MAE относится к его собственным pointwise выборам; плохие ответы не удалены.
\input{results/score_table.tex}
\input{results/outcome.tex}

\section*{Проверка posterior и производных}
На первых 32 фиксированных source-train holdout точках вычислен posterior
непрерывного исходного закона. Exp задаётся равномерными осевой координатой
и углом при радиусе $3e^{-u}$; Spiral ---
$(\sin(t^2)/t,\cos(t^2)/t)$ при равномерном $t\in[1,100]$.
Используется точный измеренный renderer исходного набора. Эталон не является
KDE по train и никогда не используется как teacher или ограничение сети.
Student интеграл вычисляется Gamma-смесью непрерывных Gaussian интегралов,
включая производную весов. Эталон учитывает все ambient координаты ядра.

В следующей таблице $\R_{\min}$ --- среднее на этих 32 точках при
$\lambda=1/256$. Последний столбец --- среднее
$|\R_{\rm model}-\R_{\rm exact}|$ на выбранном масштабе, по тем же точкам.
У изображения model здесь использует точный полный trace, чтобы убрать
Monte Carlo ошибку из сравнения с эталоном. Расхождение точного posterior
Spiral с LID 1 на конечном масштабе и расхождение сети с posterior --- разные ошибки.
\input{results/oracle_table.tex}

Дополнительно 128 test объектов зашумлены восемь раз каждый собственным
нативным шумом метода. Измерен риск
$\mathbb{E}\|b(X+\lambda Z,\lambda)-X\|^2/\lambda^2$, с суммой по всем ambient
координатам. Обученный posterior и $b_0$ получают одни и те же наблюдения и
масштаб обученной модели. Меньше лучше; отрицательный интервал парной разности
``обученный минус начальный'' подтверждает улучшение. Между ядрами наблюдения
различаются, поэтому сравнение рисков разных методов не изолирует эффект ядра.
\input{results/denoising_table.tex}

Для image test дополнительно повторяются 64 пробы с двумя независимыми seeds,
при тех же весах и уже выбранном масштабе. Масштаб не перенастраивается.
Полный trace и центральные конечные разности по всем координатам проверяются
на финальных весах в float64; отдельно проверяется float32 точность.
В следующей таблице Test MAE range --- диапазон трёх реализаций image trace
при фиксированном выбранном масштабе; у точного векторного trace повторов нет.
Предпоследний столбец --- среднее абсолютное расхождение benchmark trace с
точным trace сети на 32 holdout точках на этом масштабе. FP32 gap --- максимум
расхождения float32 и float64 trace на трёх проверочных точках при минимальном
масштабе, выбранном масштабе и $\lambda=1$. Это не ошибка относительно LID.
\input{results/probe_table.tex}

\section*{Воспроизводится ли само распределение}
Дополнительная проверка генерации заявлена после промежуточного просмотра
примерно 60000 шагов и до готовности финальных весов. Она не меняет обучение
или выбор $\lambda$. Из 512 нативных Student шумов при $\lambda=64$ получены
сэмплы интегрированием $dy/d\log\lambda=y-b(y,\lambda)$: 512 шагов Heun до
$1/256$ и итоговый posterior на этом масштабе. Это каноническая запись обоих
нативных полей. Начало с конечного масштаба приближает шумовую маргиналь;
никакой проекции на чистую геометрию в sampler нет.

Variance ratio --- отношение следов выборочных ковариаций 512 генераций и
всех 1000 test объектов; значение около 1 соответствует общему разбросу данных,
но само по себе не доказывает правильную форму распределения. Normal ratio ---
RMS расстояния сэмпла до известной линейной оболочки генератора, делённый на
RMS длину центрированного test-вектора. Эталон равен нулю. Оболочка применяется
только для этой метрики, никогда для вывода модели или траектории.
Fit distance --- среднее расстояние от generation либо test до одного
банка 8192 fit объектов, в тех же единицах; test служит контролем конечной
плотности банка. Во всех этих расстояниях сохранены полные ambient координаты.

ODE refinement --- RMS разности сэмплов на первых 64 одинаковых шумах при
512 и 1024 шагах, на координату в fit-нормированных единицах. Заданный
численный порог равен .005; он проверяет интегрирование, а не качество генерации.
Исходные 512-step сэмплы сохраняются независимо от результата этой проверки.
\input{results/generation_table.tex}
\input{results/assessment.tex}

\section*{Воспроизведение}
Producing model/config source: \texttt{b90a2b7}, ветка
\texttt{feat/non-gaussian-ambient-20260911}. Нативные реализации, unrestricted
маршруты, краткие обучения/загрузки всех 13 интерфейсов и измерения покрыты
262 тестами. Четыре настоящих 200-step preflight проверили обе новые модели
на полном векторе и полном изображении. Они не входят в benchmark результаты.
Сохранены веса, исходные split hashes, полные pointwise кривые, выборы масштаба,
независимые probe-повторы и источники каждой численной проверки.

Команды: \texttt{diagnostics/non\_gaussian\_ambient\_20260911/README.md}.
Числа и проверки: \texttt{results/metrics.csv}, \texttt{results/verification.json}.
Эта работа измеряет восемь обучений; техническая доступность матрицы
$39\times13=507$ не означает её полного прогона.

Первоисточники:
\href{https://arxiv.org/html/2410.14171v2}{t-Flow, Appendix B, Eq.164--166};
\href{https://proceedings.mlr.press/v202/xu23m/xu23m.pdf}{PFGM++, Eq.4--6};
\href{https://github.com/Newbeeer/pfgmpp/blob/main/training/loss.py}{официальный EDM loss}.

\clearpage
\includegraphics[width=\linewidth]{results/mean_response_logx.pdf}
Средние по одним 32 holdout точкам. Сплошные линии --- обученный posterior
с точным полным trace, пунктир --- непрерывный эталон. Точки отмечают масштаб,
выбранный по всем 1000 holdout объектам. Здесь показаны также малые шумы.

\clearpage
\includegraphics[width=\linewidth]{results/mean_response_linearx.pdf}
Те же неизменённые значения на линейной шкале $\lambda$.

\clearpage
\includegraphics[width=\linewidth,height=.9\textheight,keepaspectratio]{results/point_response_linearx.pdf}
Одна фиксированная holdout точка в каждой задаче/представлении;
оба метода используют один и тот же ID, указанный на панели.
\clearpage
\includegraphics[width=\linewidth,height=.92\textheight,keepaspectratio]{results/generation_coordinates.pdf}
Координаты генератора служат только для отображения сэмплов. Полная ошибка
в нормальных направлениях напечатана на каждой панели и учтена в таблице.
\clearpage
\includegraphics[width=\linewidth]{results/generation_images.pdf}
Первые четыре test изображения и первые четыре независимые генерации каждого
метода. Диапазон яркости фиксирован по всему test набору отдельно для Exp и
Spiral; нормировки контраста каждого изображения нет. Численные метрики
используют исходные значения пикселей без обрезания.
\end{document}
